\section{The Substrate and the Benchmark}
\label{sec:benchmark}

\paragraph{The substrate the literature shares.}
\label{sec:substrate}
The EarningsCall (EC) benchmark of \citet{qin2019what} and the two MAEC cohorts
are scored on split and label files redistributed by VolTAGE and KeFVP. We
fetched them, manifested them by SHA-256 and audited them
(Table~\ref{tab:substrate}; \texttt{ecvol audit substrate}). The VolTAGE and
KeFVP test files are byte-identical: one physical label set, never recomputed
from prices. Every train/validation and validation/test boundary has a zero-day
gap (on EC, training ends and validation starts on 2017-08-03) against targets
that run 30 days past each call. On the temporal splits the field reports,
80.4\% (EC), 44.2\% (MAEC-15) and 55.7\% (MAEC-16) of test calls come from
companies also in training; a temporal split always permits this, but no paper in
the lineage reports the disjoint condition, so the confound of \citet{scss2024}
has never been measured on these benchmarks. The source paper defines its horizons
over \emph{calendar} days; most downstream work describes them as trading days,
which at $\tau{=}30$ is 21 versus 30 sessions. Finally, the shipped MAEC labels
correlate only 0.52--0.87 with calendar-day targets recomputed from prices under
the same call-date anchor (0.25--0.66 at $\tau{=}3$), so the lineage's labels are
not reproducible from prices under any single documented rule.

% Curated from data/results/result_table_5r.csv and result_table_5r_labels.csv
% (runs 20261006T202254Z / 20261006T202821Z-audit-substrate-48b9432a). Overlap = test calls whose
% ticker appears in train; embargo = calendar days between the last train (val)
% date and the first val (test) date. Ours = committed splits, same metrics.
\begin{table}[t]
\centering
\small
\begin{tabular}{lrrrrr}
\toprule
Benchmark / split & calls (tr/va/te) & overlap & embargo tr$\to$va & va$\to$te & label defect \\
\midrule
EC (Qin \& Yang)         & 392 / 56 / 112  & 80.4\% & 0 d & 0 d & 137 / 16{,}800 zero cells \\
MAEC-15 (KeFVP files)    & 535 / 76 / 154  & 44.2\% & 0 d & 0 d & 4 binary cols at $\tau$ \\
MAEC-16 (KeFVP files)    & 980 / 140 / 280 & 55.7\% & 0 d & 3 d & 4 binary cols at $\tau$ \\
\midrule
ours: FinCall temporal        & 1489 / 228 / 483 & 91.3\% & 47 d & 47 d & --- \\
ours: FinCall ticker-disjoint & 1706 / 243 / 484 & \textbf{0.0\%} & n/a & n/a & --- \\
ours: MAEC temporal           & 1666 / 69 / 529 & 79.4\% & 101 d & 48 d & --- \\
ours: MAEC ticker-disjoint    & 1884 / 271 / 533 & \textbf{0.0\%} & n/a & n/a & --- \\
ours: Earnings25 ticker-disjoint & 333 / 47 / 95 & \textbf{0.0\%} & n/a & n/a & --- \\
\bottomrule
\end{tabular}
\caption{\textbf{Substrate audit of the shared benchmark files} (Result Table 5R,
regenerated by \texttt{ecvol audit substrate} from SHA-256-manifested copies of the
VolTAGE and KeFVP repositories). The VolTAGE and KeFVP split files are
byte-identical. Overlap is a property every temporal split permits, ours included;
the point is that no paper in the lineage reports the disjoint condition. The
zero-day embargos leak 30-day target windows across every boundary. Shipped
MAEC labels join our calendar-day, call-date-anchored targets on 77--79\% of test calls at
$\tau\ge7$ and correlate 0.52--0.87 with them (0.25--0.66 at $\tau{=}3$).}
\label{tab:substrate}
\end{table}


\paragraph{Corpora and identity.}
\textbf{FinCall-Surprise} \citep{fincall2025} (primary; Apache-2.0; 2{,}688 calls,
2019--2021, MP3 audio, transcripts, slides) ships four fields per record and no
ticker, date, CIK or speaker names. We reconstruct identity and date from transcript
content and slide metadata, cross-check against prices, and release the table with
an 80-row manual override list: 2{,}496 calls resolve (92.9\%; 95.3\% of
earnings-type calls) and a 60-call manual audit passes 60/60. \textbf{MAEC}
\citep{li2020maec} (secondary; CC-BY-SA-4.0; $\sim$3{,}443 calls, 2015--2018,
sentence-aligned audio features, no raw audio) and \textbf{Earnings25}
\citep{jiang2026earnings25} (CC-BY-4.0; 514 S\&P~500 calls held July--December
2025, webcast MP3 at 64\,kbps for 281 calls and $\le$24\,kbps for 232) complete the
set; Earnings25 is post-cutoff for every model we use, and 475 of its calls resolve
to a constituent ticker and join prices.

\paragraph{Targets.}
Daily adjusted prices come from a free source cross-checked against Tiingo on 50
tickers. With $r_t$ the return on
session $t$ relative to day~0, the last session before the call's information is
public, the primary target is
$v_{\mathrm{post}}(\tau)=\ln\sqrt{\frac{1}{n}\sum_{t=1}^{n}(r_t-\bar r)^2}$ for
$\tau\in\{3,7,15,30\}$, with $v_{\mathrm{pre}}(\tau)$ defined over the sessions
ending at day~0. Under the \textbf{trading-day} convention $n=\tau$; under the
\textbf{calendar-day} convention of \citet{qin2019what} the window holds every
session within $\tau$ calendar days, so $n\approx0.7\tau$ (FinCall:
2.4/4.9/10.8/21.2 sessions) and at $\tau{=}3$ at most one session for a Thursday or
Friday call. Windows with fewer than two sessions are excluded with a reason code
(42.5\% of FinCall's calls at $\tau{=}3$ calendar days, for which the released label
sets still carry a value). The conventions are far from interchangeable:
$v_{\mathrm{post}}$ correlates 0.78 between them at $\tau{=}3$ and the persistence
MSE, the leaderboard's metric, is 44--63\% higher there under calendar days, so a
published MSE is uninformative unless its convention is stated (Appendix
Table~\ref{tab:convention}). Both ship as separate files; every result below uses
trading days, and the convention is a column of every result table. Two
identity-robust variants are added: the \textbf{change}
$\Delta v=v_{\mathrm{post}}-v_{\mathrm{pre}}$, which subtracts the company's own
level, and the \textbf{HAR-residual}. After joining, FinCall yields 2{,}433 calls
with a target (97.5\%) and MAEC 2{,}688 (78.6\%; 2015--2018 delistings).

\paragraph{Measured call timing.}
\label{sec:timing}
Fewer than 4\% of FinCall transcripts mention a clock time, so no call carries a
measured time of day. We take the information boundary from EDGAR: the acceptance
time of the earliest 8-K with Item~2.02 within three days before to one day after
the call date (89.3\% of FinCall, 68.6\% of MAEC, 99.8\% of Earnings25), with
recorded fallback tiers (Earnings25 call-start time, the before/after-market flag
of \citet{scss2024}, an after-hours assumption). A release before the open makes
the call date itself day~1. Against the
call-date anchor that every released label set uses, \textbf{day~0 moves for 53\%
of FinCall, 56\% of MAEC and 72\% of Earnings25 calls}, almost always one session
earlier: under the published anchor the reaction session sits in the \emph{pre}
window for most of the corpus. On the measured anchor the persistence MSE is up to
11\% higher on FinCall, 25\% on MAEC and 27\% on Earnings25, and because the
reaction session now enters the post window, $\Delta v$ acquires a positive mean
(FinCall $+0.52$ at $\tau{=}3$, $+0.10$ at $\tau{=}30$). The call-date targets are
shipped as a labelled variant for comparison with prior label sets.

\paragraph{Splits and release.}
\emph{Temporal} splits ($\sim$70/10/20, calendar order) carry a 30-trading-day
embargo so that no target window crosses a boundary; \emph{\tickerdisjoint} splits
(grouped, seeded) put no validation or test company in training (overlap 0.0\%,
against 91.3\% FinCall and 79.4\% MAEC on our own temporal splits); a
\emph{combined} split enforces both. Every feature function receives an
\texttt{as\_of} timestamp and the harness fails any feature that reads later price
data; leakage assertions run in CI. The release contains every artefact above, the
derived-feature archives, the harness, and every result CSV with its run manifest
(resolved config, git SHA, seeds, file hash); raw audio and transcripts are never
redistributed.
